%% ma: L10-B14-0093
%% lop: 10
%% chuong: 5
%% bai: 14
%% dang: 10.14.12
%% loai: DS
%% muc: [TH, NB, VD, TH]
%% dap_an: "SĐSĐ"
%% nguon: "chuyen-de-cac-so-dac-trung-cua-mau-so-lieu-khong-ghep-nhom-toan-10-chuong-trinh-moi.pdf (GV. Phan Nhật Linh) (P05), Bài 03 Dạng 2 Phần II Câu 10"
%% kien_thuc: "Lớp 10 Bài 13 (số trung bình, trung vị, tứ phân vị, mốt), Bài 14 (khoảng biến thiên, khoảng tứ phân vị, phương sai, độ lệch chuẩn, giá trị bất thường)"
%% trang_thai: chinh-thuc
%% ghi_chu: "Tự duyệt theo yêu cầu giáo viên 22/10/2026."
\begin{ex}
Điểm một bài kiểm tra môn Toán của một nhóm học sinh cho ở bảng sau:
\begin{center}
\begin{tabular}{|c|c|c|c|c|c|c|}\hline
Điểm & $5$ & $6$ & $7$ & $8$ & $9$ & $10$\\\hline
Số học sinh & $2$ & $5$ & $6$ & $8$ & $3$ & $1$\\\hline
\end{tabular}
\end{center}
Xét tính đúng sai của các khẳng định sau:
\choiceTF{Trung vị của mẫu số liệu là $6$}{\True Khoảng biến thiên của mẫu số liệu là $R=5$}{Độ lệch chuẩn điểm của các học sinh nhỏ hơn $1$ điểm}{\True Khoảng tứ phân vị của mẫu số liệu là $\Delta_Q=2$}
\loigiai{$n=25$; tần số tích lũy $2$; $7$; $13$; $21$; $24$; $25$.

a) Sai: $M_e=x_{13}=7$.

b) Đúng: $R=10-5=5$.

c) Sai: $\overline x=7{,}32$; $s^2=1{,}5776$; $s\approx1{,}256>1$.

d) Đúng: $Q_1=\dfrac{x_6+x_7}2=6$; $Q_3=\dfrac{x_{19}+x_{20}}2=8$; $\Delta_Q=2$.}
\end{ex}
